% compendium/execution-trading-costs/execution-quality.tex
% Entries migrated on 2026-10-15 from the legacy transaction-cost-analysis file (slice 13), owned by
% market-structure-execution-quality; dispositions and source checks in notes/migration.md, the blocks
% verbatim in notes/artefacts/market-making-lob/derivations.tex.

\section{Best execution and broker conflicts}

\begin{qfdefinition}[Best execution: three standards, the dimensions of execution quality, and why retail clients see commissions but not prices]\label{def:best-execution}
\qffield{Statement} A broker taking an order assumes an agency responsibility to obtain best execution, which is not well defined. To unsophisticated clients it means the best possible price for a market order and the fastest execution for a limit order, an absolute standard, which is what the term generally means in U.S. equity markets; to more sophisticated traders, the execution they pay for; to the most sophisticated, the execution they can expect given what they pay and how well they can audit it \citep[chap.~``Brokers'', sec.~7.4.2, pp.~160--161]{alpha-harris}. Price is one dimension of execution quality and speed another; market-order traders care most about speed, limit-order traders about price, and best execution for standing limit orders depends on whether and when they execute \citep[chap.~``Internalization, Preferencing, and Crossing'', sec.~25.1, p.~515]{alpha-harris}. Kissell defines best execution as the process of choosing the strategy most likely to achieve the fund's investment objective, judged on the information available when trading began, not by closeness to a benchmark price \citep[chap.~3, p.~88]{s7-kissellscience}.

\qffield{Conditions} In U.S. equities as Harris describes them in 2002, dealers claim best execution when they fill marketable orders at the NBBO, a standard that applies only to orders smaller than the displayed size; some markets also expect average price improvement for small market orders (p.~515). Brokers that internalise or accept payment for order flow have a conflict of interest and must regularly and rigorously examine execution quality; SEC Rules 11Ac1-5 and 11Ac1-6 required market centres to publish monthly execution-quality statistics and brokers quarterly routing reports (pp.~517--518, boxes). Retail traders can audit commissions and response times easily but not price quality, which needs trade and quote data and norms for comparison; so brokers compete on commissions, and in competitive markets the total of spread and commission for small market orders does not depend on how best execution is defined (pp.~517--518). Lehalle and Laruelle contrast two regimes as of 2018: in the U.S. the trade-through rule, they write, makes venues reroute orders marketable at a better price elsewhere, which needs a consolidated pre-trade view; in Europe the executing broker chooses how to consolidate access to liquidity and must disclose and justify a best-execution policy, raising the cost of fragmentation for intermediaries \citep[sec.~3.4.2, p.~234]{alpha-laruelle}.

\qffield{Practice} Retail execution quality is measured by price against the NBBO, the midpoint or the effective spread (\hyperref[def:spread-measures]{spread measures}), speed, whether and when limit orders execute, and commissions; audit firms compare each night's trades with quotes and report exceptions (Harris, pp.~515--519). Interpretation: these measures depend on order type and size and on the quotes at the time; comparisons between venues need like orders. \emph{Best execution after fragmentation} (slice~35). In Lehalle and Laruelle's account MiFID left each broker to define best execution in its own policy, and best execution focused on ``immediate price improvement'' for aggressive orders from snapshots of visible books; after fragmentation they argue it must cover consolidated access for passive and aggressive orders, with adverse selection, passive splits and reversal as measures judged over the order originator's horizon, and must account for clearing and settlement costs that differ by venue \citep[Introduction, pp.~17--19]{alpha-laruelle}; in Europe, best-execution policies took the place of the U.S.\ trade-through rule \citep[chap.~1, sec.~1.1.2, p.~43]{alpha-laruelle}. The legacy principle (fees, fill probability, latency, hidden liquidity, access, size and horizon besides displayed price) summarises this.

\qffield{Pitfalls} Three obligations should not be merged: \emph{quote protection}, \emph{best execution} and \emph{execution-quality reporting}. Quote protection in the U.S. is the order protection rule, 17 CFR 242.611, checked on 15 October 2026 against the eCFR point-in-time text for 13 October 2026 (adopted 70 FR 37620, 29 June 2005; last amended 89 FR 26616, 15 April 2024, whose effective and compliance dates were not checked): it requires a \emph{trading center} to maintain written policies reasonably designed to prevent trade-throughs of protected quotations in NMS stocks, and to surveil them, with exceptions including system failure, non-regular-way trades, single-priced opening, reopening and closing transactions, crossed protected quotations, intermarket sweep orders (executed, or routed simultaneously against the full displayed size of protected quotations), orders whose price is not based on the quote, venues whose best quote within the previous second was equal to or worse than the trade-through price, and stopped customer orders (paragraphs (a)--(b)(9)). It is a duty of the venue, not the broker's best-execution duty, and it does not oblige a venue to reroute: an intermarket sweep order lets the router take the protected quotations itself. Lehalle and Laruelle's claims that U.S. clearing and settlement fees must be equal across venues, and their description of the European regime, were not checked against current law; the European statements here are as of their 2018 printing. The legacy principle (slice 14) that ``U.S. routing leans on protected consolidated quotes'' while ``European brokers choose the consolidation layer'' follows their p.~234. The legacy principle called best execution ``the best auditable outcome given the client's objective, limit, urgency, available venues and information risk''; that merges Harris's third standard (auditability) with Kissell's process definition, which do not coincide. Its retail ``adverse selection: post-fill price movement'' measure is a realised-spread idea (\hyperref[def:spread-measures]{spread measures}), not in the best-execution passages checked. Rules 11Ac1-5 and 11Ac1-6 are now Rules 605 and 606; the expanded Rule 605 reports required from August 2026, and the 2025 withdrawal of the SEC's best-execution proposal, are the subject of this pathway's readings \texttt{ex-sec-rule605-2024} and \texttt{ex-sec-withdrawal-2025}, not read: the regulatory statements here are as of 2002.

\qffield{See also} \hyperref[def:broker-agency-conflicts]{broker conflicts}; \hyperref[meth:transaction-cost-analysis]{transaction cost analysis}; \hyperref[def:spread-measures]{spread measures}.

\qffield{Further reading} The legacy blocks it replaces, verbatim: \nolinkurl{notes/artefacts/market-making-lob/derivations.tex}.

\qffield{Index} best execution; execution quality; NBBO; price improvement; speed of execution; limit order execution; retail execution quality; Rule 11Ac1-5; Rule 11Ac1-6; Rule 605; Rule 606; payment for order flow; internalisation; auditability; transaction audit; order protection rule; Rule 611; trade-through; protected quotation; intermarket sweep order; best-execution policy

\qffield{Status} Verification: partly checked 2026-10-15 (against the text layers of local copies of \texttt{alpha-harris}, 2002 Oxford printing, pp.~160--163 and 514--519, \texttt{s7-kissellscience}, Academic Press printing with copyright 2014, chap.~3, p.~88, and \texttt{alpha-laruelle}, 2018 World Scientific printing, p.~234; 17 CFR 242.611 checked on 15 October 2026 against the eCFR point-in-time text for 13 October 2026, slice 14; slice~35: \texttt{alpha-laruelle}, pp.~17--19 and~43). Stability: draft. Currency: Rule 611 as verified on 15 October 2026; other U.S. rules and practice as of 2002, European regime as of 2018; superseded in part, see Pitfalls.
\end{qfdefinition}

\begin{qfdefinition}[Broker conflicts: the principal--agent problem in order handling, from payment for order flow to front running]\label{def:broker-agency-conflicts}
\qffield{Statement} Brokers are agents whose interests can differ from their clients'; clients manage the problem with rewards and penalties, which require measuring performance, and measuring execution quality is hard \citep[chap.~``Brokers'', secs.~7.4--7.4.1, pp.~159--160]{alpha-harris}. Harris's conflicts (same chap.):
\begin{itemize}
\item \emph{dual trading}: a broker-dealer that internalises a client order wants the opposite price to the client, and may trade ahead of clients on the same side (sec.~7.4.3, p.~161);
\item \emph{order preferencing}: routing to a preferred dealer, an own dealer subsidiary, or in exchange for reciprocal flow or payment for order flow, about 1 cent a share in U.S. stocks in 2002, independently of current quotes (sec.~7.4.4, pp.~161--162);
\item \emph{soft dollars and directed brokerage}: commissions that buy research or other services, and sponsors directing orders to particular brokers or recapturing commissions (pp.~153--155);
\item \emph{dishonest practices}: front running, inappropriate order exposure, fraudulent trade assignment of the best fills to favoured clients, prearranged trading and kickbacks, unauthorised trading and churning (sec.~7.5, pp.~163--167).
\end{itemize}
Kissell adds venue rebates: when brokers keep maker--taker rebates or pay the fees, routing may serve the broker's net fees rather than the client \citep[chap.~3, p.~90]{s7-kissellscience}. Interpretation (the legacy proposition, slice 14): if routing to $v$ costs the client $C_v$ in expectation and earns the broker $h_v$, the client's choice is $\arg\min_vC_v$ and a broker maximising its own net would choose $\arg\min_v(C_v-h_v)$; they differ exactly when some $a,b$ have $C_a<C_b$ but $C_a-h_a>C_b-h_b$. This states when a conflict can arise; it is not evidence that brokers route that way, since reputation, regulation and competition for flow also enter the broker's objective.

\qffield{Conditions} Payments for order flow resemble kickbacks but generally are not, since competing dealers pass the value of order flow to brokers and brokers to clients through lower commissions (Harris, pp.~166 and~516--517). Dishonesty is likeliest where it is hard to detect; audit trails, supervision and reputation deter it (pp.~163--165).

\qffield{Practice} Interpretation: transaction cost analysis and audit trails turn suspected conflicts into testable comparisons, such as fill quality by venue against the NBBO or price improvement by routing destination, but a comparison shows a conflict's effect only when orders are comparable. \emph{Order-flow arrangements} (slice~35). The legacy taxonomy of internalisation, preferencing, payment for order flow and internal crossing, and its broker-agency caveat, duplicate this entry. Harris adds that brokers who cross client orders internally to earn two commissions and print them at a market with a thin book bypass standing limit orders with time precedence \citep[chap.~``Competition Within and Among Markets'', box ``Lost Time Precedence'', p.~537]{alpha-harris}. Since the 2024 MiFIR amendments, EU investment firms may not receive payment for forwarding or executing retail clients' orders, subject to a possible national exemption until 30 June 2026 (\hyperref[rule:eu-trading-venue-transparency]{EU trading venues and transparency}).

\qffield{Pitfalls} The legacy caveat on order preferencing (slice 14) restates Harris's preferencing for service, economics or reciprocal flow; its conclusion, that review must compare all-in fills rather than routing intent, is Interpretation consistent with his best-execution standards. The legacy taxonomy grouped ``allocation'' separately and listed ``commission recapture'' as an activity incentive; Harris treats recapture as a sponsor's tool to reduce the cost of soft-dollar commissions, not a broker's incentive. Figures and rules are as of 2002.

\qffield{See also} \hyperref[def:best-execution]{best execution}; \hyperref[def:order-anticipation]{order anticipation} (front running); \hyperref[def:order-exposure-decisions]{order exposure decisions}.

\qffield{Further reading} The legacy block it replaces, verbatim: \nolinkurl{notes/artefacts/market-making-lob/derivations.tex}.

\qffield{Index} principal--agent problem; broker conflicts; dual trading; internalisation; order preferencing; payment for order flow; reciprocal order flow; soft dollars; directed brokerage; commission recapture; front running; inappropriate order exposure; fraudulent trade assignment; prearranged trading; kickbacks; churning; maker--taker rebates; audit trail; routing conflict; internal crossing; payment for order flow ban (EU)

\qffield{Status} Verification: partly checked 2026-10-15 (against the text layers of local copies of \texttt{alpha-harris}, 2002 Oxford printing, pp.~153--155, 159--167 and 516--517, and \texttt{s7-kissellscience}, Academic Press printing with copyright 2014, chap.~3, p.~90; slice~35: p.~537). Stability: draft. Currency: practices and figures as of 2002 and 2014.
\end{qfdefinition}
